% SPDX-FileCopyrightText: Copyright © 2026 Idiap Research Institute <contact@idiap.ch>
% SPDX-FileContributor: Andre Freitas <andre.freitas@idiap.ch>
% SPDX-License-Identifier: MIT
% Neuro-symbolic AI Group
% File Documentation
%
% Motivation:
%   Protect Mathformer behavior while the framework evolves across Lean, ATP, KB, SymPy, docs, and skills.
%
% Intent:
%   Encode regression checks that catch parser breakage, semantic drift, and missing workflow evidence.
%
% Function:
%   Contains executable tests, fixtures, and manifests used for local and CI validation.
%
% Maintenance Notes:
%   Keep this header synchronized with major responsibility changes, and prefer updating tests or examples when behavior changes.

% Origin annotation:
% The theorem statement in this test file is sourced from:
% https://en.wikipedia.org/wiki/Sazonov%27s_theorem
\documentclass[11pt]{article}
\usepackage[T1]{fontenc}
\usepackage[utf8]{inputenc}
\usepackage{lmodern}
\usepackage{amsmath,amssymb,amsthm}
\usepackage{hyperref}

\newtheorem{theorem}{Theorem}
\newtheorem{remark}{Remark}

\title{Sazonov's Theorem}
\author{}
\date{}

\begin{document}

\maketitle

\section*{Origin Annotation}
This theorem statement was prepared as a test artifact.
Source origin: Wikipedia, \emph{Sazonov's theorem},
\url{https://en.wikipedia.org/wiki/Sazonov\%27s_theorem}.

\section*{Statement}

\begin{theorem}[Sazonov]
Let $G$ and $H$ be Hilbert spaces, and let
\[
T : G \to H
\]
be a bounded linear operator.

Assume that $T$ is a Hilbert--Schmidt operator, i.e.\ there exists an
orthonormal basis $\{e_i\}_{i\in I}$ of $G$ such that
\[
\sum_{i\in I} \|T e_i\|_H^2 < \infty.
\]

Then $T$ is $\gamma$-radonifying; that is, the pushforward under $T$ of the
canonical Gaussian cylinder set measure on $G$ extends to a bona fide
Gaussian measure on $H$.
\end{theorem}

\section*{Remark}

In the Hilbert-space setting, the converse also holds: a bounded linear
operator between Hilbert spaces is $\gamma$-radonifying if and only if it is
Hilbert--Schmidt.

\section*{References}

\begin{thebibliography}{9}

\bibitem{Schwartz1973}
L. Schwartz,
\textit{Radon Measures on Arbitrary Topological Spaces and Cylindrical Measures},
Oxford University Press, 1973.

\bibitem{WikiSazonov}
\textit{Sazonov's theorem},
Wikipedia,
\url{https://en.wikipedia.org/wiki/Sazonov\%27s_theorem}.

\end{thebibliography}

\end{document}
